\documentclass[11pt,a4paper]{article}
\usepackage[utf8]{inputenc}
\usepackage[T1]{fontenc}
\usepackage[a4paper,margin=1in]{geometry}
\usepackage{hyperref}
\usepackage{longtable}
\usepackage{booktabs}
\usepackage{array}
\usepackage{enumitem}
\usepackage{xcolor}
\usepackage{listings}

\hypersetup{
  colorlinks=true,
  linkcolor=blue,
  urlcolor=blue,
  pdftitle={Exploration Report: Agent Structure of CMBAgent and OpenGauss},
  pdfauthor={OpenAI Codex}
}

\lstdefinestyle{plaincode}{
  basicstyle=\ttfamily\small,
  frame=single,
  breaklines=true,
  columns=fullflexible
}

\title{Exploration Report:\\Agent Structure of CMBAgent and OpenGauss}
\author{Prepared for the Science-Work-Flow project}
\date{April 22, 2026}

\begin{document}
\maketitle

\begin{abstract}
This report studies the internal agent structure of the locally cloned
\texttt{cmbagent} and \texttt{OpenGauss} repositories without relying on any
remote API execution. The purpose is to understand how each system is built,
what kind of ``agent'' each project actually implements, and how the two
projects fit the evaluation framework described in
\texttt{Cambagent\_lab\_proposal.pdf}. The main result is that CMBAgent is a
role-specialized multi-agent workflow graph for scientific tasks, while
OpenGauss is a project-centered orchestration shell for Lean workflows and
managed tool-using backend sessions. They are related, but they solve different
layers of the problem.
\end{abstract}

\section{Goal and Method}
The exploration was done by reading the local source code under:
\begin{itemize}[leftmargin=1.5em]
\item \texttt{external/cmbagent}
\item \texttt{external/OpenGauss}
\end{itemize}

The focus was not on installation, credentials, or live provider access. The
focus was on structure:
\begin{itemize}[leftmargin=1.5em]
\item where the main runtime entry points are,
\item how agents are created,
\item how control and handoff are represented,
\item how tools and execution are attached,
\item how each repository stores project state, and
\item how the two systems relate to the proposal's evaluation logic.
\end{itemize}

The most informative files for CMBAgent were:
\begin{itemize}[leftmargin=1.5em]
\item \texttt{cmbagent/cmbagent.py}
\item \path{cmbagent/base_agent.py}
\item \path{cmbagent/hand_offs.py}
\item \texttt{cmbagent/functions/registration.py}
\item \path{cmbagent/workflows/one_shot.py}
\end{itemize}

The most informative files for OpenGauss were:
\begin{itemize}[leftmargin=1.5em]
\item \texttt{cli.py}
\item \path{run_agent.py}
\item \path{gauss_cli/project.py}
\item \path{gauss_cli/autoformalize.py}
\item \path{gauss_cli/commands.py}
\item \path{swarm_manager.py}
\item \texttt{toolsets.py}
\end{itemize}

\section{CMBAgent: A Scientific Multi-Agent Workflow Graph}

\subsection{Core Design Idea}
CMBAgent is built around a strong idea: scientific work is decomposed into
named agent roles, and those roles are connected by explicit handoff logic.
This means the architecture is not just ``call one model with tools.'' Instead,
it is a graph of cooperating agents such as planner, researcher, engineer,
controller, executor, summarizer, and formatter agents.

The central runtime lives in \texttt{cmbagent/cmbagent.py}. Its
\texttt{CMBAgent} class prepares the working directory, shared context,
selected model configuration, and imported agent definitions. It then registers
functions and handoffs before starting a multi-agent chat pattern.

\subsection{How Agents Are Defined}
One of the clearest structural choices in CMBAgent is that agents can be
defined in two ways:
\begin{enumerate}[leftmargin=1.5em]
\item a Python class plus YAML configuration, or
\item YAML only, using \texttt{BaseAgent} as a wrapper.
\end{enumerate}

This happens in the \texttt{import\_agents()} function in
\texttt{cmbagent.py}. The loader walks the \texttt{agents/} directory, detects
whether an agent folder has a Python file, and dynamically builds the right
class. That gives CMBAgent a hybrid architecture:
\begin{itemize}[leftmargin=1.5em]
\item simple agents stay declarative,
\item complex agents can add custom Python behavior,
\item the overall agent inventory can grow without rewriting the runtime.
\end{itemize}

\subsection{Role Inventory}
The agent tree shows that CMBAgent groups roles by function rather than by a
single monolithic prompt. The main categories are:
\begin{itemize}[leftmargin=1.5em]
\item \textbf{planning}: planner, plan\_setter, plan\_recorder, plan\_reviewer,
reviewer formatter, review recorder
\item \textbf{research}: researcher, researcher formatter, researcher executor,
summarizer, summarizer formatter
\item \textbf{coding}: engineer, engineer\_nest, engineer formatter, executor,
executor\_bash, executor formatter
\item \textbf{control}: controller, control starter, terminator
\item \textbf{hypothesis}: idea maker, idea hater, idea saver, formatter roles
\item \textbf{specialized}: camb context and camb response formatter
\item \textbf{keywords/admin/install}: keyword finders, installer, admin
\end{itemize}

This is important for the proposal because it matches an evaluation style based
on failure localization. If the planner is weak, that is a planning failure. If
the engineer writes wrong code, that is an execution or reasoning failure. If
the controller routes badly, that is a control failure. The architecture
already exposes these categories.

\subsection{Runtime Objects and Control Flow}
CMBAgent does not rely on one custom scheduler from scratch. It leans on AG2
and uses:
\begin{itemize}[leftmargin=1.5em]
\item \texttt{ContextVariables}
\item \texttt{AutoPattern}
\item group chat management
\item conditional handoff targets
\item nested chats
\end{itemize}

At a high level, the control flow is:
\begin{lstlisting}[style=plaincode]
task
-> CMBAgent(...)
-> import_agents()
-> register_functions_to_agents(...)
-> register_all_hand_offs(...)
-> initiate_group_chat(...)
-> controller / researcher / engineer / terminator loops
\end{lstlisting}

This means the system is graph-structured first and prompt-structured second.
The routing graph is part of the architecture, not a hidden side effect.

\subsection{Handoffs Are a First-Class Abstraction}
The file \path{hand_offs.py} is one of the most revealing files in the whole
repository. It explicitly wires agent-to-agent flow.

Examples of the structural patterns in that file:
\begin{itemize}[leftmargin=1.5em]
\item simple linear chains such as planner $\rightarrow$ formatter
$\rightarrow$ recorder
\item conditional routing based on context variables, such as whether feedback
rounds remain
\item mode-dependent routing, for example when response formatters return to the
controller versus a direct engineer path
\item message-history reduction for formatter and recorder roles
\item nested chat construction for the engineer execution loop
\item LLM-driven controller branching with choices such as engineer,
researcher, idea-maker, idea-hater, or terminator
\end{itemize}

This gives CMBAgent a recognizable workflow-engine flavor. It is not merely
``multi-agent'' in a loose marketing sense. It is a controlled graph with
defined transitions and explicit branch points.

\subsection{Why the Controller Matters}
The controller is the main orchestration brain once the graph is live. In
\path{hand_offs.py}, the controller can route work to different role agents
based on language-model conditions. That is a powerful but also risky design:
\begin{itemize}[leftmargin=1.5em]
\item it makes the workflow adaptive,
\item but it also creates a specific failure mode where incorrect routing can
propagate downstream errors.
\end{itemize}

For the proposal, this is valuable rather than inconvenient. A core goal of the
proposal is to characterize failure modes of scientific agents, and routing
mistakes are a concrete class of such failures.

\subsection{Execution Model}
CMBAgent has a strong separation between reasoning agents and execution agents.
The engineer proposes code-oriented work; executors run or save outputs; format
agents compress results; the controller decides what happens next. This
separation is useful for scientific evaluation because it lets the benchmark ask
where an error came from:
\begin{itemize}[leftmargin=1.5em]
\item bad planning,
\item bad scientific reasoning,
\item bad code generation,
\item bad tool execution,
\item bad summarization,
\item or bad routing.
\end{itemize}

\subsection{Architectural Strengths of CMBAgent}
\begin{itemize}[leftmargin=1.5em]
\item Very explicit role decomposition
\item Good match for workflow-level evaluation
\item Natural support for failure taxonomies
\item Hybrid YAML plus Python agent definitions keep the inventory extensible
\item Clear place for domain-specific scientific context agents
\end{itemize}

\subsection{Architectural Limitations of CMBAgent}
\begin{itemize}[leftmargin=1.5em]
\item The routing logic is powerful, but complexity grows quickly
\item Behavior depends on several interacting agent prompts and control edges
\item Some failures are emergent across handoffs rather than local to one file
\item External provider assumptions are still present in default model paths
\end{itemize}

\section{OpenGauss: A Project-Centered Lean Workflow Orchestrator}

\subsection{Core Design Idea}
OpenGauss is structurally different. It is not primarily a static graph of many
named in-process roles. It is a project shell around Lean work. Its main job is
to help a user or backend agent operate inside a Lean project, invoke workflows
such as \texttt{/prove} and \texttt{/formalize}, and manage child sessions,
tooling, and project metadata.

The most important architectural decision is that OpenGauss treats the
\textbf{project} as the stable object, not the individual agent role.

\subsection{The Project Model Is Foundational}
The file \path{gauss_cli/project.py} shows that OpenGauss begins by finding
or validating a Lean project. It recognizes a project root, expects a
\texttt{.gauss/project.yaml} manifest, validates schema details, locates the
Lean root, and creates runtime/cache/workflow directories.

That means OpenGauss stores state in a much more project-native way than
CMBAgent. The important object is \texttt{GaussProject}, which packages:
\begin{itemize}[leftmargin=1.5em]
\item project root
\item manifest path
\item Lean root
\item runtime/cache/workflow directories
\item source metadata
\item blueprint markers
\end{itemize}

This makes OpenGauss feel closer to a developer toolchain or research shell than
to a pure benchmark agent.

\subsection{User Interface and Command Surface}
The top-level \texttt{cli.py} and \path{gauss_cli/commands.py} make it clear
that OpenGauss is meant to be operated interactively. Its user model is:
\begin{enumerate}[leftmargin=1.5em]
\item start the Gauss CLI in or for a project,
\item issue slash commands,
\item let the system launch or manage workflows,
\item inspect, attach to, or continue sessions.
\end{enumerate}

This is a strong contrast with CMBAgent's more programmatic workflow entry
points such as \texttt{one\_shot} and \texttt{deep\_research}.

\subsection{The Generic Tool-Calling Agent}
The file \path{run_agent.py} contains the reusable runtime class
\texttt{AIAgent}. This class is the closest OpenGauss equivalent to a general
agent engine. It manages:
\begin{itemize}[leftmargin=1.5em]
\item model/provider selection
\item tool calling loops
\item iteration budgets
\item context compression
\item prompt assembly
\item scratchpad and trajectory handling
\item tool execution and error recovery
\end{itemize}

This is a broad, reusable runtime, but it is still not the same architecture as
CMBAgent. Here, the emphasis is on a robust single agent runtime with tools and
managed child sessions, not on a fixed many-role graph.

\subsection{Managed Workflow Launching}
The file \path{gauss_cli/autoformalize.py} reveals OpenGauss's real
specialization. It normalizes workflow commands such as:
\begin{itemize}[leftmargin=1.5em]
\item \texttt{/prove}
\item \texttt{/draft}
\item \texttt{/review}
\item \texttt{/checkpoint}
\item \texttt{/refactor}
\item \texttt{/golf}
\item \texttt{/autoprove}
\item \texttt{/formalize}
\item \texttt{/autoformalize}
\end{itemize}

It then builds managed launch plans, backend-specific runtime context, staged
assets, and handoff requests for child sessions. In other words, OpenGauss
specializes in \emph{workflow dispatch over Lean projects}.

This explains why OpenGauss is a strong complement to the proposal, but not a
drop-in substitute for CMBAgent. It is excellent for formalization and theorem
workflow support, but it is not designed around cosmology-style multi-agent
scientific planning and execution.

\subsection{Swarm Management and Sessions}
The file \path{swarm_manager.py} adds another important layer: background
workflow management. OpenGauss tracks task identifiers, progress, process state,
recent output, attach/detach behavior, and per-task Lean status. This makes it
possible to treat long-running formalization work as managed sessions rather
than as one opaque model call.

That is architecturally significant because it means OpenGauss is built for:
\begin{itemize}[leftmargin=1.5em]
\item persistent work on projects,
\item long-lived background jobs,
\item reattachment,
\item and developer-style control over sessions.
\end{itemize}

\subsection{Tool Discipline}
The file \texttt{toolsets.py} shows another deliberate design choice. Gauss
keeps its default tool surface narrow:
\begin{itemize}[leftmargin=1.5em]
\item file operations
\item web search
\item browser actions
\end{itemize}

This is a smaller and more disciplined tool surface than one might expect from
a broad autonomous agent framework. The advantage is that it narrows the
behavioral space around the actual tasks OpenGauss wants to support.

\subsection{Architectural Strengths of OpenGauss}
\begin{itemize}[leftmargin=1.5em]
\item Strong project model for Lean work
\item Explicit managed workflows for proving and formalization
\item Good session and swarm control for long-running tasks
\item Clear CLI entry points that users can operate directly
\item Good fit for complementary evaluation on formalization tasks
\end{itemize}

\subsection{Architectural Limitations of OpenGauss}
\begin{itemize}[leftmargin=1.5em]
\item It is not primarily a scientific multi-agent benchmark graph
\item Domain reasoning is shaped by the Lean project and workflow layer
\item Evaluation needs to focus on formalization outcomes, not just generic chat
\item Its architecture is less naturally decomposed into planner/researcher/
engineer roles than CMBAgent
\end{itemize}

\section{Direct Comparison}

\begin{longtable}{>{\raggedright\arraybackslash}p{0.24\textwidth} >{\raggedright\arraybackslash}p{0.34\textwidth} >{\raggedright\arraybackslash}p{0.34\textwidth}}
\toprule
\textbf{Dimension} & \textbf{CMBAgent} & \textbf{OpenGauss} \\
\midrule
\endfirsthead
\toprule
\textbf{Dimension} & \textbf{CMBAgent} & \textbf{OpenGauss} \\
\midrule
\endhead
\bottomrule
\endfoot
Primary unit of organization &
Named role agents connected by handoffs &
Lean project plus workflow commands and managed sessions \\

Main runtime style &
Graph-like multi-agent orchestration &
Interactive CLI plus backend launch orchestration \\

Core control logic &
Controller routes among planner, researcher, engineer, idea, and termination roles &
Commands create launch plans; managed backends execute workflows inside project context \\

State model &
Shared context variables, message history transforms, working directory &
\texttt{GaussProject}, runtime/cache/workflow directories, child session metadata \\

Agent identity &
Many specialized agents with explicit task roles &
Generic tool-calling agent runtime plus workflow-specific launch behavior \\

Execution style &
Scientific reasoning, coding, formatting, summarizing, control loops &
Formalization, proof, review, refactor, and project-scoped Lean assistance \\

Best evaluation fit &
Scientific workflow reliability and failure taxonomy &
Formalization support and Lean workflow quality \\

Main weakness for the proposal &
Complex control graph can hide emergent failure chains &
Less aligned with cosmology workflow structure as a primary agent benchmark \\
\end{longtable}

\section{Implications for the Proposal}

\subsection{Why CMBAgent Should Stay Central}
The proposal in \texttt{Cambagent\_lab\_proposal.pdf} is mainly about evaluating
scientific agent reliability across planning, tool-grounded work, inference,
and error propagation. CMBAgent maps naturally onto that because:
\begin{itemize}[leftmargin=1.5em]
\item its internal roles already reflect those stages,
\item its controller makes routing errors observable,
\item its executor split makes reasoning versus execution easier to separate,
\item its domain-specific context agents allow science-specific grounding.
\end{itemize}

\subsection{Why OpenGauss Should Be Complementary}
OpenGauss should remain in the project, but as a complementary component. It is
especially useful for:
\begin{itemize}[leftmargin=1.5em]
\item formalizing mathematical claims that appear in scientific workflows,
\item testing whether informal reasoning can be translated into Lean-style
checked structure,
\item comparing free-form scientific reasoning with stricter formal artifacts,
\item and running experiments on a real Lean project such as
\texttt{projects/LeanCourse25}.
\end{itemize}

In proposal terms, OpenGauss is best viewed as \textbf{support for the
evaluation framework}, not as the same kind of evaluated agent as CMBAgent.

\subsection{Recommended Minimal Division of Labor}
\begin{itemize}[leftmargin=1.5em]
\item Use CMBAgent as the primary object of scientific workflow evaluation.
\item Use OpenGauss when the benchmark needs a formalization or proof-oriented
companion task.
\item Compare them only where the task definition makes that comparison fair.
\item Avoid judging OpenGauss as though it were a cosmology workflow graph, and
avoid judging CMBAgent as though it were a dedicated Lean project shell.
\end{itemize}

\section{A Practical Mental Model}
If the project is viewed as a lab, then the simplest mental model is:
\begin{itemize}[leftmargin=1.5em]
\item \textbf{CMBAgent} is the experimental scientific workflow system being
studied.
\item \textbf{OpenGauss} is a specialized companion environment for formal and
Lean-related tasks that can support parts of the evaluation program.
\end{itemize}

This single distinction removes much of the confusion between the two
repositories. They both use LLM-driven agents, but they are architected for
different layers of work.

\section{Conclusion}
The deep structural difference between the repositories is now clear.
CMBAgent is a multi-agent workflow graph with explicit role decomposition and
handoff control. OpenGauss is a project-centered orchestration environment for
Lean workflows, with a generic tool-calling agent runtime and managed child
sessions.

For the proposal, this means CMBAgent should remain the main evaluated agent
framework, while OpenGauss should be kept minimal and complementary. That
division is not a compromise; it is the cleanest architectural reading of the
two codebases.

\end{document}
